% GENERATED FILE. Edit canonical lessons, then run npm run generate:books.
% canonical-source-hash: fnv1a64:e117b29593b80ca6
% canonical-lessons: ES-C514-reportaje, ES-C514-reportero, ES-C514-entrevistar, ES-C514-difundir, ES-C514-anunciar

\chapter{El reportaje}
\label{ch:es-el-reportaje}
\hlchaptermodality{514}

\begin{chapteropening}
\textbf{By the end of this chapter:} \emph{I can say that a reporter interviewed someone, spread the news and announced something.}

Complete the last lesson of chapter 514: I can say that a reporter interviewed someone, spread the news and announced something.
\end{chapteropening}

\section[el reportaje]{el reportaje — the report, the feature story}

\label{lesson:ES-C514-reportaje}



\begin{quote}
Before the new one: say the Spanish for the correspondent, then the Spanish for to publish.
\end{quote}

\subsection*{You'll want to know: el reportaje}
\textbf{el reportaje} — \textquotedblleft{}the report, the feature story\textquotedblright{}.

\begin{quote}
Vimos un reportaje sobre la ciudad.
\end{quote}

\begin{quote}
\emph{We watched a report about the city.}
\end{quote}

A longer piece in a paper or on television. A short official report is \emph{un informe}.

\begin{cousinweb}
\emph{reportaje} was borrowed from French \emph{reportage}, which is usually traced to \emph{reporter}, from Latin \emph{reportāre} 'to carry back' (Lewis \& Short). A report carries the story back to you.
\end{cousinweb}

\begin{grammarlens}[title={What you've built}]
One more word for telling, explaining and describing at B1.
\end{grammarlens}

\subsection*{Your turn}
\begin{itemize}
\raggedright
  \item \emph{Say it:} \emph{el reportaje}
  \item \emph{Say it:} \emph{el reportaje}, once more
  \item \emph{Say it:} a sentence of your own with \emph{el reportaje}
  \item \emph{Recall:} say the Spanish for the correspondent, then the Spanish for to publish, then say \emph{el reportaje} again
\end{itemize}

\begin{tcolorbox}[breakable,colback=teal!4,colframe=teal!35!black,title={Before you move on}]
What does el reportaje mean? (\textquotedblleft{}The report, the feature story\textquotedblright{}.) Say it once more.
\end{tcolorbox}
\section[el reportero]{el reportero — the reporter}

\label{lesson:ES-C514-reportero}



\begin{quote}
Before the new one: say the Spanish for to publish, then the Spanish for the report.
\end{quote}

\subsection*{You'll want to know: el reportero}
\textbf{el reportero} — \textquotedblleft{}the reporter\textquotedblright{}.

\begin{quote}
El reportero llegó antes que la policía.
\end{quote}

\begin{quote}
\emph{The reporter arrived before the police.}
\end{quote}

A woman is \emph{la reportera}.

\begin{cousinweb}
\emph{reportero} is usually traced to English \emph{reporter}, given the Spanish ending \emph{-ero}. The English word comes through French from Latin \emph{reportāre} 'to carry back' (Lewis \& Short).
\end{cousinweb}

\begin{grammarlens}[title={What you've built}]
One more word for telling, explaining and describing at B1.
\end{grammarlens}

\subsection*{Your turn}
\begin{itemize}
\raggedright
  \item \emph{Say it:} \emph{el reportero}
  \item \emph{Say it:} \emph{el reportero}, once more
  \item \emph{Say it:} a sentence of your own with \emph{el reportero}
  \item \emph{Recall:} say the Spanish for to publish, then the Spanish for the report, then say \emph{el reportero} again
\end{itemize}

\begin{tcolorbox}[breakable,colback=teal!4,colframe=teal!35!black,title={Before you move on}]
What does el reportero mean? (\textquotedblleft{}The reporter\textquotedblright{}.) Say it once more.
\end{tcolorbox}
\section[entrevistar]{entrevistar — to interview}

\label{lesson:ES-C514-entrevistar}



\begin{quote}
Before the new one: say the Spanish for the report, then the Spanish for the reporter.
\end{quote}

\subsection*{You'll want to know: entrevistar}
\textbf{entrevistar} — \textquotedblleft{}to interview\textquotedblright{}.

\begin{quote}
Entrevistaron al alcalde en la radio.
\end{quote}

\begin{quote}
\emph{They interviewed the mayor on the radio.}
\end{quote}

\emph{La entrevista} is the interview. It takes personal \emph{a} with people: \emph{entrevistar a alguien}.

\begin{cousinweb}
\emph{entrevistar} is built on \emph{entrevista}, made from \emph{entre} 'between' and \emph{vista} 'a sight', and usually said to follow French \emph{entrevue}. \emph{Vista} comes from Latin \emph{vidēre} 'to see' (EtymDB).
\end{cousinweb}

\begin{grammarlens}[title={What you've built}]
One more word for telling, explaining and describing at B1.
\end{grammarlens}

\subsection*{Your turn}
\begin{itemize}
\raggedright
  \item \emph{Say it:} \emph{entrevistar}
  \item \emph{Say it:} \emph{entrevistar}, once more
  \item \emph{Say it:} a sentence of your own with \emph{entrevistar}
  \item \emph{Recall:} say the Spanish for the report, then the Spanish for the reporter, then say \emph{entrevistar} again
\end{itemize}

\begin{tcolorbox}[breakable,colback=teal!4,colframe=teal!35!black,title={Before you move on}]
What does entrevistar mean? (\textquotedblleft{}To interview\textquotedblright{}.) Say it once more.
\end{tcolorbox}
\section[difundir]{difundir — to spread, to broadcast}

\label{lesson:ES-C514-difundir}



\begin{quote}
Before the new one: say the Spanish for the reporter, then the Spanish for to interview.
\end{quote}

\subsection*{You'll want to know: difundir}
\textbf{difundir} — \textquotedblleft{}to spread, to broadcast\textquotedblright{}.

\begin{quote}
La radio difundió la noticia en pocos minutos.
\end{quote}

\begin{quote}
\emph{The radio spread the news in a few minutes.}
\end{quote}

Used for news, ideas and images: \emph{difundir una foto}.

\begin{cousinweb}
\emph{difundir} is a learned borrowing from Latin \emph{diffundere} 'to pour out, to spread' (EtymDB), built on \emph{dis-} 'apart' and \emph{fundere} 'to pour'.
\end{cousinweb}

\begin{grammarlens}[title={What you've built}]
One more word for telling, explaining and describing at B1.
\end{grammarlens}

\subsection*{Your turn}
\begin{itemize}
\raggedright
  \item \emph{Say it:} \emph{difundir}
  \item \emph{Say it:} \emph{difundir}, once more
  \item \emph{Say it:} a sentence of your own with \emph{difundir}
  \item \emph{Recall:} say the Spanish for the reporter, then the Spanish for to interview, then say \emph{difundir} again
\end{itemize}

\begin{tcolorbox}[breakable,colback=teal!4,colframe=teal!35!black,title={Before you move on}]
What does difundir mean? (\textquotedblleft{}To spread, to broadcast\textquotedblright{}.) Say it once more.
\end{tcolorbox}
\section[anunciar]{anunciar — to announce}

\label{lesson:ES-C514-anunciar}



\begin{quote}
Before the new one: say the Spanish for to interview, then the Spanish for to spread.
\end{quote}

\subsection*{You'll want to know: anunciar}
\textbf{anunciar} — \textquotedblleft{}to announce\textquotedblright{}.

\begin{quote}
El presidente anunció su viaje a Chile.
\end{quote}

\begin{quote}
\emph{The president announced his trip to Chile.}
\end{quote}

\emph{El anuncio} is the announcement, and also an advert.

\begin{cousinweb}
\emph{anunciar} is a learned borrowing from Latin \emph{annūntiāre} 'to make known' (EtymDB), built on \emph{nūntius} 'a messenger'. English \emph{announce} is the same Latin verb through French.
\end{cousinweb}

\begin{grammarlens}[title={What you've built}]
One more word for telling, explaining and describing at B1.
\end{grammarlens}

\subsection*{Your turn}
\begin{itemize}
\raggedright
  \item \emph{Say it:} \emph{anunciar}
  \item \emph{Say it:} \emph{anunciar}, once more
  \item \emph{Say it:} a sentence of your own with \emph{anunciar}
  \item \emph{Recall:} say the Spanish for to interview, then the Spanish for to spread, then say \emph{anunciar} again
\end{itemize}

\begin{tcolorbox}[breakable,colback=teal!4,colframe=teal!35!black,title={Before you move on}]
What does anunciar mean? (\textquotedblleft{}To announce\textquotedblright{}.) Say it once more.
\end{tcolorbox}
